% !TEX root = ../main.tex
\chapter{Tensor Storage Model}
\label{ch:tensor_model}

A tensor framework's storage model is the type definition of its
central data structure --- in \Lucid's case, what fields the
\code{TensorImpl} class carries, how those fields relate to backing
memory, and what invariants the engine guarantees between them. The
storage model is invisible to the user, who manipulates a Python
\code{Tensor} that hides everything below, but it determines the
performance and correctness envelope of every operation. A storage
model that does not support views forces every reshape to copy. A
storage model that does not separate stride from shape forces every
permutation to materialise. A storage model whose backing buffer is
owned by the wrong layer cannot be shared zero-copy with the GPU.

This chapter formalises \Lucid's storage model. The model is
deliberately conventional --- it is the same model that any
NumPy-derived framework uses, with one important difference: the
backing storage may belong to \MLX\ rather than to \Lucid's own
allocator, which makes the bridge of \chref{ch:unified_memory}
possible without breaking the abstraction. The conventional pieces
are still worth describing in their own right, because the rest of
the book assumes them.

\secref{sec:tensor_model:tensorimpl} describes the \code{TensorImpl}
class and its fields. \secref{sec:tensor_model:storage} describes
the \code{Storage} class and its ownership semantics.
\secref{sec:tensor_model:views} describes view operations and what
makes them cheap. \secref{sec:tensor_model:autograd_meta} describes
the autograd metadata that hangs off each tensor.
\secref{sec:tensor_model:refcounts} describes the reference
counting that ties the pieces together.
\secref{sec:tensor_model:bridge} describes how \MLX\ arrays enter
the storage model. \secref{sec:tensor_model:invariants} states the
invariants the rest of the engine relies on.

\section{The TensorImpl Class}
\label{sec:tensor_model:tensorimpl}

\code{TensorImpl} is defined in
\code{lucid/\_C/tensor/TensorMeta.h}. It is a C++ struct, not an
abstract base class --- there is exactly one concrete
\code{TensorImpl} type, and the differences between CPU tensors and
GPU tensors live in the \code{Storage} that \code{TensorImpl}
points to, not in subclassing.
\figref{fig:tensor_model:layout} sketches the object graph: a
Python \code{Tensor} holds a \code{TensorImpl}, which holds shape /
stride / dtype / device / offset metadata plus shared pointers to a
\code{Storage} (the backing buffer) and an \code{AutogradMeta} (the
optional differentiation state).

\begin{figure}[t]
\centering
\begin{adjustbox}{max width=\textwidth, center}
\begin{tikzpicture}[
  pyobj/.style={draw, rounded corners=4pt, font=\small\sffamily,
                fill=lucid.note.bg, minimum width=2.8cm,
                minimum height=1.2cm, align=center},
  impl/.style={draw, rounded corners=2pt, font=\small\sffamily,
               fill=lucid.info.bg, minimum width=4.5cm,
               minimum height=2.6cm, align=left, inner sep=4pt},
  storage/.style={draw, rounded corners=2pt, font=\small\sffamily,
                  fill=lucid.ok.bg, minimum width=5.2cm,
                  minimum height=2.4cm, align=left, inner sep=4pt},
  meta/.style={draw, rounded corners=2pt, font=\small\sffamily,
               fill=lucid.code.bg, minimum width=4.0cm,
               minimum height=2.0cm, align=left, inner sep=4pt},
  ptr/.style={->, >=stealth, thick, color=lucid.accent.dark,
              shorten >=2pt, shorten <=2pt},
  shared/.style={->, >=stealth, dashed, color=lucid.muted,
                 shorten >=2pt, shorten <=2pt},
  ptrlbl/.style={font=\footnotesize\ttfamily, fill=white,
                 inner sep=2pt},
]
  \node[pyobj] (pytensor) at (-6.5, 2) {\textbf{Tensor}\\(Python)};
  \node[impl] (impl) at (0, 2) {
    \textbf{TensorImpl}\\[2pt]
    \texttt{shape}, \texttt{strides}\\
    \texttt{dtype}, \texttt{device}\\
    \texttt{storage\_offset}\\
    \texttt{version\_counter}
  };
  \node[storage] (storage) at (8, 4) {
    \textbf{Storage}\\[2pt]
    \texttt{data}: \texttt{void*}\\
    \texttt{nbytes}\\
    \texttt{provider}: \\
    \quad Lucid/Mlx/External\\
    \texttt{deleter}: callback
  };
  \node[meta] (autograd) at (8, 0) {
    \textbf{AutogradMeta}\\[2pt]
    \texttt{requires\_grad}\\
    \texttt{grad\_fn}\\
    \texttt{grad}, \texttt{is\_leaf}\\
    (null if unused)
  };

  \draw[ptr] (pytensor) -- node[ptrlbl, above]
    {shared\_ptr} (impl);
  \draw[ptr] (impl.north east) -- node[ptrlbl, above, sloped]
    {shared\_ptr} (storage.south west);
  \draw[ptr] (impl.south east) -- node[ptrlbl, below, sloped]
    {unique\_ptr} (autograd.north west);

  % Shared storage: another tensor pointing at same Storage (view)
  \node[impl, dashed, opacity=0.85, minimum height=2cm]
       (impl2) at (0, -2.5) {
    \textbf{TensorImpl (view)}\\[2pt]
    different \texttt{shape},\\
    different \texttt{strides},\\
    same \texttt{storage} ptr
  };
  \draw[shared] (impl2.east) to[out=0, in=-90]
    node[ptrlbl, midway, below right] {share storage}
    (storage.south);
\end{tikzpicture}
\end{adjustbox}
\caption[Tensor object graph.]{The \code{TensorImpl} object graph.
The Python \code{Tensor} wraps a \code{shared\_ptr<TensorImpl>}.
The \code{TensorImpl} holds metadata inline and points (via
\code{shared\_ptr}) to a heap-allocated \code{Storage} that owns
the backing buffer; multiple \code{TensorImpl}s may share the same
\code{Storage} (view aliasing, shown dashed). The
\code{AutogradMeta} pointer is null for non-differentiable tensors,
which keeps the per-tensor overhead low for the majority of
intermediates.}
\label{fig:tensor_model:layout}
\end{figure}

The conceptual layout is:

\begin{lstlisting}[style=lucidcpp,language=C++]
struct TensorImpl {
  Shape                          shape;
  Strides                        strides;
  Dtype                          dtype;
  Device                         device;
  size_t                         storage_offset;
  std::shared_ptr<Storage>       storage;
  std::unique_ptr<AutogradMeta>  autograd_meta;  // null if !requires_grad
  uint64_t                       version_counter;
};
\end{lstlisting}

Each field is small and explicit. Together they fully describe a
tensor.

\subsection{Shape}
\label{ssec:tensor_model:shape}

\code{shape} is a small-vector of dimension sizes. The element type
is a 64-bit signed integer to permit individual dimensions larger
than \(2^{31}\) (rare but possible for streaming workloads), and
the vector itself uses small-buffer optimisation: ranks up to 6 are
stored inline without heap allocation, ranks 7 and above promote to
a heap allocation. Most practical tensors have rank at most 5
(batch, channels, height, width, sometimes a depth), so the inline
storage is hit by the overwhelming majority of allocations.

A zero-rank tensor (a scalar) has empty shape and a single element.
A zero-size tensor has at least one dimension equal to zero and no
elements at all; the storage may still be non-null but is unused.

\subsection{Strides}
\label{ssec:tensor_model:strides}

\code{strides} is a small-vector of the same length as \code{shape}.
Each entry is the number of elements (not bytes) to step in the
backing storage to advance by one along that dimension. The element
type is signed because negative strides are permitted (they encode
reversed views).

For a contiguous tensor of shape \((d_0, d_1, \ldots, d_{n-1})\) in
C-order (the default), the strides are
\((d_1 d_2 \cdots d_{n-1}, d_2 \cdots d_{n-1}, \ldots, d_{n-1}, 1)\).
For a Fortran-ordered tensor, the strides are reversed. \Lucid\
uses C-order by default; the dispatcher does not assume contiguity
but does prefer it.

\subsection{Dtype and device}
\label{ssec:tensor_model:dtype_device}

\code{dtype} is an enum identifying the element type. The enum is
defined in \code{core/Dtype.h} and includes the entries discussed
in \chref{ch:mlx_runtime}~Section~\ref{ssec:mlx_runtime:array_dtype}:
\code{float16}, \code{float32}, \code{float64}, \code{bfloat16},
\code{int8} through \code{int64} (signed only), \code{bool}, and
\code{complex64}.

\code{device} is similarly an enum from \code{core/Device.h} with
exactly two values: \code{cpu} and \code{metal}. The values are
exhaustive on Apple silicon: \Lucid\ does not expose
\code{cuda}, \code{xla}, or any other device label.

\subsection{Storage offset}
\label{ssec:tensor_model:storage_offset}

\code{storage\_offset} is the number of elements from the start of
the backing \code{Storage} to the first element of this tensor. For
a freshly allocated tensor, \code{storage\_offset} is zero. For a
sliced view, \code{storage\_offset} encodes the position of the
slice's first element.

The offset is in elements, not bytes, to keep dtype-aware indexing
simple. The byte offset is \code{storage\_offset * sizeof(dtype)}.

\subsection{The Storage handle}
\label{ssec:tensor_model:storage_handle}

\code{storage} is a \code{shared\_ptr} to a reference-counted
\code{Storage} object. Multiple \code{TensorImpl} instances may
share the same \code{Storage} (this is exactly what makes views
zero-copy), and the \code{Storage}'s underlying buffer is freed
only when the last referring \code{TensorImpl} is destroyed.

\subsection{The autograd metadata pointer}
\label{ssec:tensor_model:autograd_ptr}

\code{autograd\_meta} is a \code{unique\_ptr} to an
\code{AutogradMeta} object, allocated only when the tensor is
created with \code{requires\_grad=True}. The pointer is null for
the great majority of tensors (intermediate activations, optimiser
internals, sliced views), which saves both memory and the cost of
walking the autograd graph for tensors that do not participate in
backpropagation. \secref{sec:tensor_model:autograd_meta} treats
\code{AutogradMeta} in detail.

\subsection{The version counter}
\label{ssec:tensor_model:version_counter}

\code{version\_counter} is a 64-bit counter that increments on
every in-place modification of the tensor. It is the mechanism by
which autograd detects ``a saved tensor for backward has been
modified in place'' --- a class of bug that would silently corrupt
gradients --- and raises an exception instead.

The counter starts at zero on creation. \code{add\_}, \code{mul\_},
and other in-place ops increment it. When the autograd engine saves
a tensor for backward, it captures the current counter value;
during backward, it checks that the counter is unchanged. A mismatch
raises an error with a message that names the in-place operation
responsible.

\section{The Storage Class}
\label{sec:tensor_model:storage}

\code{Storage} is the heap-allocated, reference-counted buffer that
multiple \code{TensorImpl}s may share. Its definition is in
\code{lucid/\_C/tensor/Storage.h}:

\begin{lstlisting}[style=lucidcpp,language=C++]
struct Storage {
  void*               data;     // raw pointer to backing buffer
  size_t              nbytes;   // total bytes
  Dtype               dtype;    // (redundant cache for fast checks)
  Device              device;   // ditto
  StorageProvider     provider; // who allocated this buffer
  std::function<void(void*)> deleter;
};
\end{lstlisting}

\subsection{Provider tag}
\label{ssec:tensor_model:provider}

The \code{provider} tag distinguishes between buffers allocated
through different paths:

\begin{itemize}
  \item \code{LucidPool}: the buffer came from \Lucid's own
    size-class pool (\chref{ch:allocator}).
  \item \code{Mlx}: the buffer is owned by an \code{mlx::core::array}.
    The \code{Storage} holds a \code{shared\_ptr<array>} in a
    type-erased slot so that the array stays alive as long as the
    \code{Storage} does.
  \item \code{Accelerate}: the buffer was allocated through
    \Accelerate's \code{vDSP\_create\_*} helpers. Rare; used by a
    small set of CPU kernels that prefer aligned scratch space.
  \item \code{External}: the buffer was provided by an external
    library (numpy, DLPack producer) and \Lucid\ does not own it.
    The deleter is the foreign release function.
\end{itemize}

The provider tag has no effect on operations --- a kernel that
reads the buffer does not care who allocated it --- but it
determines how the buffer is freed. The \code{deleter} closure is
constructed at \code{Storage} creation time and called when the
\code{Storage}'s reference count drops to zero.

\subsection{Reference counting}
\label{ssec:tensor_model:refcount}

\code{Storage} is heap-allocated and managed through
\code{shared\_ptr}. The reference count is the standard
\code{shared\_ptr} count: every \code{TensorImpl} pointing at the
\code{Storage} contributes one. When the last \code{TensorImpl} is
destroyed, the \code{Storage}'s \code{shared\_ptr} count reaches
zero, the deleter is invoked, and the underlying memory is freed.

The pattern is the C++ idiom for shared ownership and requires no
special discipline from the engine. The only subtlety is that
\code{TensorImpl}'s autograd graph (the chain of \code{grad\_fn}
nodes) may hold references to intermediate \code{Storage}s that the
user no longer has handles to; this is why memory usage grows during
a forward pass and shrinks after a backward, rather than tracking
the Python live set directly.

\section{Views}
\label{sec:tensor_model:views}

A \emph{view} is a \code{TensorImpl} whose \code{storage} pointer
is shared with another \code{TensorImpl} (its source) and whose
shape, strides, or \code{storage\_offset} differ. Views are the
mechanism by which \Lucid\ supports reshape, transpose, slice, and
expand without copying data.

\subsection{What operations produce views}
\label{ssec:tensor_model:view_ops}

The following operations return views of their input, without
copying:

\begin{itemize}
  \item \code{reshape(shape)}: if the new shape is compatible with
    the source's strides (i.e., the source is contiguous in the
    relevant axes), returns a view. Otherwise materialises.
  \item \code{transpose(dim0, dim1)}: always a view. Strides are
    permuted; shape is permuted.
  \item \code{permute(dims)}: always a view. Generalised transpose.
  \item \code{slice(dim, start, end, step)}: always a view. Shape
    is adjusted; \code{storage\_offset} advances; stride may scale
    by \code{step}.
  \item \code{select(dim, index)}: always a view. Rank decreases
    by one.
  \item \code{narrow(dim, start, length)}: always a view.
    Equivalent to \code{slice} with step 1.
  \item \code{expand(shape)}: always a view. New dimensions of
    size larger than 1 are represented by stride 0; the underlying
    data is not duplicated.
  \item \code{squeeze}, \code{unsqueeze}: always views.
  \item \code{flip(dim)}: always a view. Strides become negative
    along the flipped dimension; \code{storage\_offset} advances
    to the end.
\end{itemize}

The contract: any of these operations runs in constant time (no
data motion). The \code{Storage} pointer is shared with the source
and the new \code{TensorImpl} has independent shape/strides/offset
fields.

\subsection{When a view is not contiguous}
\label{ssec:tensor_model:noncontig}

A view's strides may or may not match the canonical contiguous
strides for its shape. If they do, the view is contiguous. If they
do not, the view is non-contiguous, and many kernels that assume
unit-stride along the innermost dimension cannot operate on it
directly.

The dispatcher handles this in one of two ways depending on the
kernel:

\begin{itemize}
  \item Kernels that natively support strided input (most
    \Accelerate\ vDSP primitives, most \MLX\ primitives) pass the
    strides through. No copy.
  \item Kernels that do not (some convolution paths, the
    SharedStorage bridge to CPU) materialise the view via a
    \code{contiguous()} call. One copy.
\end{itemize}

The cost of the contiguous-materialisation is the cost of one
read-write pass over the tensor's elements. It is amortised by the
fact that most operations downstream of a transpose or slice will
re-materialise once and then reuse; the policy is to delay
materialisation as long as it is permitted.

\subsection{The \texttt{is\_contiguous} predicate}
\label{ssec:tensor_model:is_contig}

Whether a tensor is contiguous is a derived property, computed from
\code{shape} and \code{strides} on demand. The check is the
standard one: walk shape and strides in reverse, verifying that
each stride equals the product of the trailing shape entries. A
zero-size tensor is trivially contiguous. A scalar is trivially
contiguous.

\code{TensorImpl::is\_contiguous()} caches the result on first
call; subsequent in-place operations that change strides invalidate
the cache. The caching matters because the predicate is checked on
many fast paths.

\section{The AutogradMeta Class}
\label{sec:tensor_model:autograd_meta}

\code{AutogradMeta} carries the per-tensor differentiation state.
It is allocated only when needed --- specifically, when a tensor is
created with \code{requires\_grad=True} or when a \code{grad\_fn}
is attached during the forward pass of a differentiable operation.

The structure:

\begin{lstlisting}[style=lucidcpp,language=C++]
struct AutogradMeta {
  bool                          requires_grad;
  std::shared_ptr<Node>         grad_fn;
  uint32_t                      output_nr;     // which output of grad_fn
  std::shared_ptr<TensorImpl>   grad;          // accumulated gradient
  std::vector<Hook>             pre_hooks;
  std::vector<Hook>             post_hooks;
  bool                          is_leaf;
};
\end{lstlisting}

\code{requires\_grad} is the user-facing flag. \code{grad\_fn} is
the backward node (\chref{ch:autograd_engine}) that this tensor's
gradient flows through. \code{grad} is the accumulated gradient
that lives on leaf tensors (parameters). \code{is\_leaf}
distinguishes parameters (created by the user) from intermediates
(produced by ops); leaves accumulate gradient on backward, while
intermediates only propagate.

\subsection{Why AutogradMeta is allocated lazily}
\label{ssec:tensor_model:lazy_alloc}

A typical training step produces tens of thousands of intermediate
\code{TensorImpl} objects, of which only a small fraction are
leaves and only the leaves need the gradient accumulator. Allocating
an \code{AutogradMeta} per tensor unconditionally would waste tens
of bytes per tensor times tens of thousands per step, which is
megabytes of garbage on a hot path.

The lazy-allocation pattern allocates the metadata only when
either (a) the user sets \code{requires\_grad=True}, or (b) the
autograd engine attaches a \code{grad\_fn} during a differentiable
operation. The null-pointer test for ``does this tensor participate
in autograd at all'' is the fast-path predicate that most code
checks.

\section{Reference Counting and Lifetime}
\label{sec:tensor_model:refcounts}

The lifetime of a tensor's storage is determined by the union of
several reference counts. Understanding the union is necessary for
predicting when memory is reclaimed.

\subsection{The user-visible Tensor count}
\label{ssec:tensor_model:user_count}

The Python \code{Tensor} object holds a \code{shared\_ptr<TensorImpl>}.
The Python reference count of the \code{Tensor} object determines
when the \code{shared\_ptr} is dropped; when it is dropped, the
\code{TensorImpl}'s reference count is decremented; if that count
reaches zero, the \code{TensorImpl}'s \code{Storage shared\_ptr} is
dropped; and so on.

The simple case is: a Python variable goes out of scope, the
\code{Tensor} is collected, the \code{TensorImpl} is freed, and the
\code{Storage} (if no other view shares it) is freed too. The
backing buffer returns to its provider's pool.

\subsection{The autograd graph count}
\label{ssec:tensor_model:autograd_count}

The complication arises when the tensor participates in autograd.
The forward pass attaches \code{grad\_fn} nodes to each output
tensor, and those nodes hold \code{shared\_ptr<TensorImpl>}
references to any tensor saved for backward. So a tensor that the
user has dropped may still be alive, retained by the autograd graph,
until \code{backward()} is called and the graph is freed.

This is why memory usage during a forward pass grows: each
operation saves the inputs (or outputs) it needs for its backward,
and those saved tensors are pinned until backward releases them.
After backward, the autograd graph is freed and all transitive
references are dropped, returning the memory to the pool.

\subsection{The view chain}
\label{ssec:tensor_model:view_chain}

Views share \code{Storage} with their source. If the user creates a
view of a large tensor, then drops the source, the source's
\code{Storage} stays alive as long as the view does. The view is a
small object (one \code{TensorImpl} with shape and strides) but it
holds a reference to the full source's backing buffer.

This can be surprising: a small slice of a large tensor keeps the
full tensor's memory alive. The user-visible escape is
\code{view.clone()}, which forces a fresh \code{Storage} with only
the view's bytes.

\section{Bridge to MLX Arrays}
\label{sec:tensor_model:bridge}

The previous sections described the conventional storage model. The
unconventional piece is the bridge to \MLX: an \code{mlx::core::array}
is the canonical storage for any GPU tensor in \Lucid, and the
\code{Storage} wraps the \MLX\ array rather than owning its own
buffer.

\subsection{How a GPU tensor's storage is constructed}
\label{ssec:tensor_model:gpu_storage}

When \Lucid\ allocates a GPU tensor, the path is:

\begin{enumerate}
  \item The dispatcher determines the device is \code{metal}.
  \item It calls \code{MlxBackend::allocate(shape, dtype)}, which
    calls \code{mlx::core::zeros(shape, dtype, mx.gpu)} (or the
    appropriate fill primitive) and returns the resulting
    \code{mlx::core::array}.
  \item It constructs a \code{Storage} with \code{provider = Mlx},
    \code{data} pointing to the array's \code{data<T>()}, and the
    deleter set to a closure that holds the \code{shared\_ptr<array>}
    alive for the storage's lifetime.
  \item It constructs a \code{TensorImpl} wrapping the
    \code{Storage} with the requested shape, strides, dtype, and
    device.
\end{enumerate}

The result: a \Lucid\ tensor whose backing storage is owned by
\MLX. Operations on this tensor go through the dispatcher to
\MLX\ primitives, which produce new \MLX\ arrays, which become new
\code{Storage}s, and so on.

\subsection{How a CPU tensor's storage is constructed}
\label{ssec:tensor_model:cpu_storage}

For a CPU tensor, the path is symmetric but simpler:

\begin{enumerate}
  \item The dispatcher determines the device is \code{cpu}.
  \item It calls \code{Allocator::alloc(nbytes)} on the size-class
    pool, getting back a raw pointer.
  \item It constructs a \code{Storage} with \code{provider = LucidPool},
    \code{data} pointing to the pool buffer, and the deleter set
    to \code{Allocator::free}.
  \item It constructs a \code{TensorImpl} wrapping the
    \code{Storage}.
\end{enumerate}

The CPU path does not involve \MLX. CPU tensors are wholly owned by
\Lucid; only the GPU path delegates ownership to \MLX.

\subsection{Round-tripping}
\label{ssec:tensor_model:round_trip}

A tensor that lives on \code{metal} and is moved to \code{cpu}
(via \code{.to('cpu')} or implicitly through a CPU op) goes through
the \code{SharedStorage} bridge of
\chref{ch:unified_memory}~Section~\ref{sec:unified_memory:bridge}.
The bridge produces a CPU pointer to the same physical bytes that
the \MLX\ array holds, wraps it in a new \code{Storage} of provider
\code{Mlx} (the storage is still owned by \MLX), and constructs a
new \code{TensorImpl} with \code{device = cpu}.

The result is two \code{TensorImpl}s, one with \code{device = metal}
and one with \code{device = cpu}, sharing the same underlying
bytes through different \code{Storage}s that both retain references
to the same \code{mlx::core::array}. This is what makes the
zero-copy bridge possible at the storage-model level: the
provider tag is not the same as the device label, and a CPU-device
tensor can have an \code{Mlx} provider.

\section{Invariants}
\label{sec:tensor_model:invariants}

We close with a short list of invariants the rest of the engine
relies on. Each is enforced either by construction or by an
assertion in \code{TensorImpl}'s constructor.

\begin{enumerate}
  \item \textbf{Shape and stride parity.} \code{shape.size() ==
    strides.size()}.
  \item \textbf{Storage non-null.} For any tensor with nonzero
    total element count, \code{storage} is non-null.
  \item \textbf{Element count fits in storage.} For any tensor,
    \code{storage\_offset + max\_addressable\_index <
    storage->nbytes / sizeof(dtype)}.
  \item \textbf{Device consistency.} \code{device ==
    storage->device}.
  \item \textbf{Dtype consistency.} \code{dtype ==
    storage->dtype} (with the carve-out that \code{storage->dtype}
    can be a byte-equivalent dtype for views that reinterpret;
    relevant only for \code{tensor.view(dtype=...)} which \Lucid\
    does not currently expose).
  \item \textbf{Version monotonicity.} \code{version\_counter} is
    monotone non-decreasing over the tensor's lifetime.
  \item \textbf{Autograd leaf consistency.} If \code{is\_leaf} is
    true, \code{grad\_fn} is null. Conversely, if \code{grad\_fn}
    is non-null, \code{is\_leaf} is false.
\end{enumerate}

A violation of any invariant produces a \code{LucidError} from the
constructor or from the asserting code path. Production builds keep
the assertions enabled; the cost is negligible.

\section{Summary and Forward References}
\label{sec:tensor_model:summary}

\Lucid's storage model is the conventional shape/stride/dtype/device
plus storage tuple, with two distinguishing features: the
\code{Storage}'s \code{provider} tag permits the underlying buffer
to be owned by \MLX\ rather than by \Lucid, which is how the
unified-memory bridge works at the type-system level; and the
\code{AutogradMeta} is allocated lazily, which keeps the per-tensor
overhead low for the great majority of intermediates.

The next chapter, \chref{ch:allocator}, treats the size-class
allocator that backs the \code{LucidPool} provider, and the memory
accounting that exposes its state to the user. \chref{ch:backend_dispatch}
treats the dispatcher that routes per-op work to the backends.
\chref{ch:bridge_boundaries} catalogues the six places where the
storage model crosses into external libraries.

\begin{lucidnote}
The user does not see the storage model directly. Through the
Python \code{Tensor} class the user sees \code{.shape},
\code{.dtype}, \code{.device}, and the autograd-related attributes
(\code{.requires\_grad}, \code{.grad}, \code{.grad\_fn}). Strides
are exposed through \code{.stride()} but are rarely consulted
directly. \code{Storage}, \code{provider}, and the bridge
mechanics are entirely internal.
\end{lucidnote}
